\documentclass{article}
\usepackage[margin=1in]{geometry}
\usepackage{siunitx}
\usepackage{amsmath}
\sisetup{per-mode=symbol,round-mode=places,round-precision=2}
\title{Specific Heat Capacity Worksheet}
\author{}
\date{}
\begin{document}
\maketitle

\section*{Questions}

\begin{enumerate}
\item 
\begin{enumerate}
    \item State the relationship between heat energy $Q$, mass $m$, specific heat capacity $c$ and temperature change $\Delta T$.
    \item Rearrange to make (i) $c$ and (ii) $m$ the subject.
    \item Give the SI units of $Q$, $m$, $c$ and $\Delta T$.
\end{enumerate}

\item A \SI{0.335}{\gram} sample of water warms from \SI{24.5}{\degreeCelsius} to \SI{26.4}{\degreeCelsius}. How much heat does the water absorb? \textit{($c_\text{water}= \SI{4180}{\joule\per\kilogram\celsius}$)}

\item How much heat must be removed from \SI{0.225}{\kilogram} of water to cool it from \SI{25.0}{\degreeCelsius} to \SI{10.0}{\degreeCelsius}? \textit{($c_\text{water}= \SI{4180}{\joule\per\kilogram\celsius}$)}

\item How much energy is needed to raise \SI{1.0}{\kilogram} of water from \SI{25}{\degreeCelsius} to \SI{99}{\degreeCelsius}?

\item An insulated cup holds \SI{0.075}{\kilogram} of water at \SI{24.00}{\degreeCelsius}. A \SI{0.026}{\kilogram} piece of metal at \SI{82.25}{\degreeCelsius} is dropped in and the mixture settles at \SI{28.34}{\degreeCelsius}. Calculate the specific heat capacity of the metal.

\item A calorimeter has a heat capacity of \SI{1265}{\joule\per\degreeCelsius}. Its temperature rises from \SI{22.34}{\degreeCelsius} to \SI{25.12}{\degreeCelsius}. How much heat was released inside the calorimeter?

\item It takes \SI{192}{\joule} to raise the temperature of \SI{45.0}{\gram} of silicon by \SI{6.0}{\degreeCelsius}. Determine the specific heat capacity of silicon.

\item One 350‑g can of cola (assume $c=\SI{4180}{\joule\per\kilogram\celsius}$) is cooled from \SI{25}{\degreeCelsius} to \SI{3}{\degreeCelsius}. How much heat is removed?

\item \SI{96}{\joule} of heat raises the temperature of a \SI{0.075}{\kilogram} block of lead by \SI{10}{\degreeCelsius}. Find the specific heat capacity of lead.

\item A \SI{0.033}{\kilogram} block of titanium warms by \SI{5.20}{\degreeCelsius} when it absorbs \SI{89.7}{\joule}. Calculate the specific heat capacity of titanium.
\end{enumerate}

\newpage
\section*{Answer Key}

\begin{enumerate}
\item 
\begin{enumerate}
    \item $Q = mc\Delta T$
    \item $c = \dfrac{Q}{m\Delta T} \quad\;\; m = \dfrac{Q}{c\Delta T}$
    \item $Q$: \si{\joule}; $m$: \si{\kilogram}; $c$: \si{\joule\per\kilogram\celsius}; $\Delta T$: \si{\celsius}
\end{enumerate}

\item $Q = (0.000335)(4180)(1.9) \approx \SI{2.7}{\joule}$

\item $Q = (0.225)(4180)(-15.0) \approx \SI{-1.41e4}{\joule}$ (heat removed)

\item $Q = (1.0)(4180)(74) \approx \SI{3.09e5}{\joule}$

\item $Q_w = (0.075)(4180)(28.34-24.00)=\SI{1.36e3}{\joule}$\\
$c_\text{metal}= \dfrac{1.36\times10^{3}}{0.026\times53.91}\approx \SI{9.7e2}{\joule\per\kilogram\celsius}$

\item $Q = (1265)(2.78) \approx \SI{3.52e3}{\joule}$ released

\item $c = \dfrac{192}{0.045\times6.0}\approx \SI{7.1e2}{\joule\per\kilogram\celsius}$

\item $Q = (0.350)(4180)(-22)=\SI{-3.22e4}{\joule}$ removed

\item $c = \dfrac{96}{0.075\times10}=\SI{1.28e2}{\joule\per\kilogram\celsius}$

\item $c = \dfrac{89.7}{0.033\times5.20}\approx \SI{5.23e2}{\joule\per\kilogram\celsius}$
\end{enumerate}

\end{document}
